% Canonical preamble for this skill. See references/latex-house-style.md.
% The look is references/house-style/ (style-spec.md, housestyle.sty); this
% file adds only the teaching affordances and the math kit on top of it.
% Replace <COURSE>, <DOC>, <NAME> and <YEAR> as needed.
% It ends before \begin{document}, so a smoke document can \input it as-is.
% Build with scripts/build.sh (lualatex, three passes); it puts housestyle.sty
% on the search path and the .sty finds the vendored fonts itself.

\documentclass[11pt]{article}
% A4: uncomment the next line; the margins stay the spec's.
% \newcommand\hspaper{a4paper}
\usepackage{housestyle}  % geometry, fonts, colour tokens, running head, the twelve blocks, hyperref (hidelinks)
\usepackage{amssymb,amsthm}
\usepackage{mathtools}
\usepackage{tabularx,multirow}

% --- Running head and foot -------------------------------------------------
% Slug at the left of the head; the section name at the right follows each
% \section through \markboth, read at shipout, so a heading pushed to the next
% page never names itself on the page before (override with \hssection{...}).
% Foot: the copyright line at the left, the page number at the right. The
% first page has no head, per the spec.
\hsslug{<COURSE> <DOC>}
\fancyfoot[L]{\hsrunlabel{\textcopyright{} <YEAR> <NAME>. All rights reserved.}}
\renewcommand{\sectionmark}[1]{\markboth{#1}{}}
\hssection{\leftmark}

% --- Headings --------------------------------------------------------------
% Headings are unnumbered on the page (the .sty's \section). Cross-reference
% a section by name with \nameref{sec:...}; a Part by its number, Part~\ref{part:...}.
% A Part opens on a new page: its eyebrow, then the page heading at 17/21 in the
% Subhead cut (the .sty's \section), no bold.
\makeatletter
\renewcommand{\thepart}{\arabic{part}}
\newcommand{\hs@parteyebrow}{}
\titleclass{\part}{top}
\assignpagestyle{\part}{fancy}  % keep the head and the copyright foot on a Part page
\titleformat{\part}[display]{\hssubheadfont\fontsize{17pt}{21pt}\selectfont\color{ink}}%
  {\hsaccentlabel{Part \thepart\hs@parteyebrow}}{10pt}{}
\titlespacing*{\part}{0pt}{0pt}{24pt}
\titleformat{\subsection}{\hssubheadfont\fontsize{13pt}{17pt}\selectfont\color{ink}}{}{0pt}{}
\titlespacing*{\subsection}{0pt}{16pt}{4pt}
\titleformat{\subsubsection}{\itshape\color{ink}}{}{0pt}{}
\titlespacing*{\subsubsection}{0pt}{12pt}{2pt}

% --- Taxonomy: category as structure, never as colour ------------------------
% \catbanner{A}{Title} opens a Part: eyebrow "Part N . Category A", then the
% heading, and the running head's section becomes the title.
\newcommand{\catbanner}[2]{%
  \gdef\hs@parteyebrow{\ \textperiodcentered\ Category #1}%
  \part{#2}\markboth{#2}{}\gdef\hs@parteyebrow{}}
\makeatother
% Numbered sections (1, 1.1, ...): uncomment. Off by default.
% \hsnumbersections
% Justified prose: uncomment for an essay, a write-up or notes that run as
% continuous prose. Leave it off for a technical document, or any document
% whose prose is broken up by figures, tables, stats and code (style-spec.md,
% Alignment).
% \hsjustified
\newcommand{\cat}[1]{\hslabel{[#1]}}      % [A]
\newcommand{\tool}[1]{\hslabel{[T#1]}}    % [T5]
\newcommand{\tools}[1]{\hslabel{[#1]}}    % [T4,T5]

% --- Teaching affordances, inside the twelve blocks -------------------------
% insight: the section's lead line, right under the heading.
\newcommand{\insight}[1]{\hslead{#1}}
% derivation: a display-equation run (block 12) the reader can skip on
% re-read. A label opens it, a grey rule shows where it ends.
\newenvironment{derivation}%
  {\par\vspace{4pt}\hslabel{Derivation \textperiodcentered\ skip on re-read}\par\nopagebreak}%
  {\par\nopagebreak\vspace{2pt}\noindent\hsthinrule\par}
% \onsheet: this equation lives on the formula sheet. Grey, not accent.
\newcommand{\onsheet}{\par\nopagebreak\noindent\hslabel{On the formula sheet}\nopagebreak}
% \opt: out of scope. Inline after a mention, or first thing in an optional paragraph.
\newcommand{\opt}{\leavevmode\unskip\hspace{4pt}\hslabel{[optional]}\hspace{4pt}\ignorespaces}
% \probhead{source}{method}{category}: one small-caps label line over a worked problem.
\newcommand{\probhead}[3]{%
  \par\vspace{8pt}\noindent\hslabel{#1 \textperiodcentered\ #2 \textperiodcentered\ #3}\par\nopagebreak}

% --- Tables ----------------------------------------------------------------
% L{w} and R{w} come from the .sty; Y is a ragged-right tabularx column.
\newcolumntype{Y}{>{\raggedright\arraybackslash}X}

% --- Math kit ----------------------------------------------------------------
\newcommand{\R}{\mathbb{R}}
\newcommand{\N}{\mathbb{N}}
\newcommand{\Z}{\mathbb{Z}}
\newcommand{\RR}[1]{\mathbb{R}^{#1}}
\newcommand{\grad}{\nabla}
\newcommand{\Hess}{H_f}

% Bold vectors -- define all lowercase so agents can grab any letter
\newcommand{\ba}{\mathbf{a}}  \newcommand{\bb}{\mathbf{b}}  \newcommand{\bc}{\mathbf{c}}
\newcommand{\bd}{\mathbf{d}}  \newcommand{\be}{\mathbf{e}}  \newcommand{\bff}{\mathbf{f}}
\newcommand{\bg}{\mathbf{g}}  \newcommand{\bh}{\mathbf{h}}  \newcommand{\bi}{\mathbf{i}}
\newcommand{\bj}{\mathbf{j}}  \newcommand{\bk}{\mathbf{k}}  \newcommand{\bl}{\mathbf{l}}
\newcommand{\bm}{\mathbf{m}}  \newcommand{\bn}{\mathbf{n}}  \newcommand{\bp}{\mathbf{p}}
\newcommand{\bq}{\mathbf{q}}  \newcommand{\br}{\mathbf{r}}  \newcommand{\bs}{\mathbf{s}}
\newcommand{\bt}{\mathbf{t}}  \newcommand{\bu}{\mathbf{u}}  \newcommand{\bv}{\mathbf{v}}
\newcommand{\bw}{\mathbf{w}}  \newcommand{\bx}{\mathbf{x}}  \newcommand{\by}{\mathbf{y}}
\newcommand{\bz}{\mathbf{z}}  \newcommand{\bzero}{\mathbf{0}}

% Transpose: braces (not ^) so X^\trans gives X^{\!\top}
\newcommand{\trans}{{\!\top}}

% Partial derivatives
\newcommand{\pderiv}[2]{\dfrac{\partial #1}{\partial #2}}
\newcommand{\ppderiv}[2]{\dfrac{\partial^2 #1}{\partial #2^2}}
\newcommand{\pmderiv}[3]{\dfrac{\partial^2 #1}{\partial #2 \, \partial #3}}

% Big-O -- use \Oh{n} or \Ohof{n \log n} everywhere; never bare $O()$
\newcommand{\Oh}{\mathcal{O}}
\newcommand{\Ohof}[1]{\mathcal{O}\!\left(#1\right)}

% Differential d
\newcommand{\diff}[1]{\mathop{}\!\mathrm{d}#1}

% --- Code --------------------------------------------------------------------
% Block 11: \hslisting{Listing 1 / what it shows} then
% \begin{Verbatim}[bgcolor=codefill] ... \end{Verbatim}. No syntax colour in print.

\hsslug{PCB-0017-A stereo-class-d-amp}
\fancyfoot[L]{\hsrunlabel{PCB-0017-A stereo-class-d-amp}}
\setlength{\tabcolsep}{4pt}
\begin{document}

\hstitleblock{PCB-0017 Rev A \textperiodcentered\ stereo-class-d-amp \textperiodcentered\ September 2026}%
{Owner's guide: stereo Class-D amplifier}%
{What the board does, how to wire and power it up, what is on it, what it
costs, and what still has to happen before anyone orders it.}

PCB-0017-A stereo-class-d-amp is a small, 2-layer stereo Class-D audio amplifier built
around the Texas Instruments TPA3118D2, run from a single 12 V supply. It
takes stereo line-level audio in on wire pads and drives two bridge-tied
speaker outputs. The board reproduces, and is compared against, a published
44 x 34 mm reference design. All four build gates -- ERC,
placement, routed DRC and DFM -- pass. The reference board used the TPA3116D2,
sold only in a pad-up package that needs a heatsink bolted on top; this
board uses the TPA3118D2 instead, whose pad-down package sinks its heat
straight into the board's own copper. That extra ground-pour area is most of
why this board, at 53.55 x 36.8 mm, runs about 32\% larger than the
reference's 44 x 34 mm.

\begin{hsplate}{top.png}{Top side, as routed}{53.55 x 36.8 mm, 2 layers}{render}%
{U1 (TPA3118D2) sits centre-board with ground pour on both sides for
heat-sinking. J1 (power) and J2 (line in) are on one edge; J3 and J4
(left/right speaker out) and TP1 (GND probe) are on the other.}
\end{hsplate}

This is a fab package, not an order: nothing has been sent to JLCPCB. Four
proposed verification waivers on U1's power and output pins (see
\nameref{sec:using}) need the owner's sign-off in
\texttt{reports/verify-waivers.proposed.json} first; the ordering steps
near the end of this guide are written for that later day.

\section{Schematic}
\hslead{One sheet: power in, line in, the amplifier, and the two LC output
filters.}

\begin{hsfigure}{Figure 1 / stereo-class-d-amp schematic}%
{Slate wire pads carry power and audio in and out; U1 is the amplifier; the
inductor-capacitor pairs on the right are the two output filters.}
\hsdiagram{../schematic.pdf}
\end{hsfigure}

Power comes in at J1 and splits to two bulk capacitors, C13 and C14, one at
each of U1's PVCC pin pairs. Line audio comes in at J2 and is AC-coupled into
U1 through C1--C4. U1 drives four switching outputs, one pair per channel,
each into its own 10~uH inductor (L1--L4) and a damping branch (R4--R7,
C23--C26), which filter the switching carrier before the signal reaches the
speaker pads J3 and J4. Every ground on the board is one net, whose
bottom-layer pour is both the current return and U1's heat spreader.

\section{Using it}
\label{sec:using}
\hslead{Four wire-pad groups, one probe pad, no switches.}

\begin{hsblock}
\hslabel{Table 1 / connectors}\par\vspace{8pt}
{\small
\begin{tabular}{L{0.08\linewidth} L{0.37\linewidth} L{0.39\linewidth}}
\hstoprule
\hshead{Pad} & \hshead{Carries} & \hsaccenthead{Notes} \\
\hline
J1 & +12V, GND & 10.8--13.2 V DC, sized for 4 A. No reverse-polarity
  protection: wire it right. \\
J2 & L, GND, R & Stereo line level, about 1 Vrms, AC-coupled. \\
J3 & Left speaker +/- & Bridge-tied load: neither pad is ground. \\
J4 & Right speaker +/- & Bridge-tied load: neither pad is ground. \\
TP1 & GND & Probe pad for a scope or meter ground lead. \\
\hline
\end{tabular}}
\end{hsblock}

Gain is fixed by R1 at 20 dB; there is no volume control, mute switch or
status LED, since R2 and R3 hold U1 always on. First power-up: connect a
speaker to J3 and/or J4 before applying power, without tying either speaker
pad to ground -- the output is bridge-tied. Apply 12 V at J1 watching
polarity, then feed line audio to J2; TP1 should read J1's GND potential.

Four checks in \texttt{reports/verify-waivers.proposed.json} are proposed,
not yet approved, on this layout: the 2-layer thermal model under U1 (the
basis instead is TI's measured 22 C/W on its 2-layer EVM; with sixteen
thermal vias under U1's pad I estimate this board at about 30 C/W, so near
85 C at the 1.5 W design point), two short
pin-pitch-limited necks on the +12V and speaker-OUT nets, and one ground-pour
check that mis-reads a deliberately split, fully-stitched fill as
disconnected. None changes the board; all four wait on the owner's name in the
\texttt{approved} field of each of that file's eight entries, one per finding.

\section{What's on it}
\hslead{39 fitted parts across 16 BOM lines: 7 unique Basic parts, 9 unique
Extended. J1--J4 and TP1 are bare pads, not JLC-assembled parts.}

\begin{hsblock}
\hslabel{Table 2 / bill of materials}\par\vspace{8pt}
{\footnotesize
\begin{tabular}{L{0.17\linewidth} L{0.26\linewidth} L{0.15\linewidth} L{0.13\linewidth} L{0.06\linewidth}}
\hstoprule
\hshead{Designators} & \hshead{Value} & \hshead{LCSC} & \hsaccenthead{Class} & \hshead{Qty} \\
\hline
C1,C2,C3,C4 & 1uF 50V X7R & C28323 & Basic & 4 \\
C5 & 1uF 25V X7R & C106858 & Extended & 1 \\
C6,C9,C10 & 100nF 50V X7R & C14663 & Basic & 3 \\
C7,C8 & 1nF 50V NP0 & C106246 & Extended & 2 \\
C11,C12 & 10uF 25V X5R & C15850 & Basic & 2 \\
C13 & 100uF 35V polymer & C46550467 & Extended & 1 \\
C14 & 100uF 35V elec. & C2992596 & Extended & 1 \\
C15,C16,C17,C18 & 220nF 50V X7R & C5378 & Basic & 4 \\
C19,C20,C21,C22 & 680nF 50V X7R & C107133 & Extended & 4 \\
C23,C24,C25,C26 & 10nF 50V X7R & C15195 & Basic & 4 \\
D1 & SMBJ15A TVS & C113988 & Extended & 1 \\
L1,L2,L3,L4 & 10uH 5.5A & C207842 & Extended & 4 \\
R1 & 5.6k 1\% & C23189 & Basic & 1 \\
R2,R3 & 100k 1\% & C25803 & Basic & 2 \\
R4,R5,R6,R7 & 3.3R 1\% & C137725 & Extended & 4 \\
U1 & TPA3118D2DAPR & C46497 & Extended & 1 \\
\hline
\end{tabular}}
\end{hsblock}

Each unique Extended part adds one of JLC's per-part feeder fees on top of
the parts themselves (see the assembly line in Table 3). J1, J2, J3, J4 and
TP1 are bare copper wire pads with no LCSC part and no connector body; hand
wiring in solder, they are not part of JLC's assembly.

\section{What it costs}
\hslead{An estimate and a lower bound, not a quote: get the real number in
JLC's own cart before ordering.}

Every figure below comes from a reference pricing table, not JLC's live API:
measured live prices have run 1.9--3.1x higher than this table. Treat every
number here as a floor, and confirm at \hsurl{https://cart.jlcpcb.com/quote}
before ordering.

\begin{hsblock}
\hslabel{Table 3 / quote matrix, HASL finish}\par\vspace{8pt}
{\small
\begin{tabular}{L{0.10\linewidth} L{0.18\linewidth} L{0.19\linewidth} L{0.17\linewidth} L{0.17\linewidth}}
\hstoprule
\hshead{Qty} & \hshead{PCB} & \hshead{Assembly} & \hshead{Board total} & \hsaccenthead{Per unit} \\
\hline
5  & \$2.00 & \$59.09 & \$61.09 & \$12.22 \\
10 & \$2.09\footnotemark[1] & \$60.18\footnotemark[1] & \$62.18 & \$6.22 \\
30 & \$8.90\footnotemark[1] & \$64.53\footnotemark[1] & \$73.43 & \$2.45 \\
\hline
\end{tabular}}\par\vspace{4pt}
{\small\color{inkseventy}\footnotemark[1]\,Split between PCB and assembly
not separately itemised at this quantity in the source data; the board
total and per-unit figures are as quoted.}
\end{hsblock}

Assembly setup and stencil fees (\$16 flat) and per-part feeder fees dominate
the cost at low quantity; the bare board itself is a couple of dollars. The
parts alone (not counting the PCB or assembly line) are estimated at
\$2.53 per board. ENIG finish prices the same as HASL in this table at every
quantity shown.

\section{Ordering at JLCPCB}
\hslead{Written for after the owner approves the four waivers above --
nothing here has been done yet.}

Before uploading anything, buy the Extended parts into JLC's inventory for
5 boards: its BOM review can otherwise show one as idle stock, unselected
at qty 0, until it has been bought in and the review re-run.

\begin{hsblock}
\hslabel{Table 4 / pre-buy list, 5 boards}\par\vspace{8pt}
{\footnotesize
\begin{tabular}{L{0.15\linewidth} L{0.26\linewidth} L{0.18\linewidth} L{0.10\linewidth} L{0.10\linewidth}}
\hstoprule
\hshead{LCSC} & \hshead{Part} & \hsaccenthead{Designators} & \hshead{Per bd} & \hshead{Buy} \\
\hline
C106858   & 1uF 25V X7R        & C5           & 1 & 5  \\
C106246   & 1nF 50V NP0        & C7,C8        & 2 & 10 \\
C46550467 & 100uF 35V polymer  & C13          & 1 & 5  \\
C2992596  & 100uF 35V elec.    & C14          & 1 & 5  \\
C107133   & 680nF 50V X7R      & C19--C22     & 4 & 20 \\
C113988   & SMBJ15A            & D1           & 1 & 5  \\
C207842   & 10uH 5.5A          & L1--L4       & 4 & 20 \\
C137725   & 3.3R 1\%           & R4--R7       & 4 & 20 \\
C46497    & TPA3118D2DAPR      & U1           & 1 & 5  \\
\hline
\end{tabular}}
\end{hsblock}

Then, in order: (1) upload
\texttt{PCB-0017-A\_stereo-class-d-amp\_gerbers.zip}; (2) check the quoted layer
count, size and thickness (2 layers, 53.55 x 36.8 mm, 1.6 mm) match the spec
above; (3) enable PCB Assembly, top side, for the build quantity; (4) upload
\texttt{BOM.csv} and \texttt{CPL.csv}; (5) confirm every part matched its
LCSC number, especially the nine Extended ones just bought in; (6) in the
placement preview, check polarity and rotation -- no automatic rotation
corrections were flagged for this board, but check by eye anyway: D1 (TVS
polarity), U1 pin 1, and the C13/C14 polarised capacitors; (7) run JLCDFM as
a second opinion; (8) pay. Payment is the owner's step, not a session's.

\section{What it cost to design}
\hslead{Model spend on generating this board -- not part of the board's own
price.}

\begin{hsblock}
\renewcommand{\arraystretch}{1.1}
\hslabel{Table 5 / generation cost by step}\par\vspace{8pt}
{\small
\begin{tabular}{L{0.55\linewidth} R{0.35\linewidth}}
\hstoprule
\hshead{Step} & \hsaccenthead{USD} \\
\hline
Before the workspace (brief and skill read) & \$0.50 \\
Brief and intake & \$1.16 \\
Research & \$3.84 \\
Architecture & \$14.22 \\
Schematic & \$3.10 \\
Placement & \$4.88 \\
Routing & \$0.43 \\
Verification & \$0.51 \\
DFM, order files and guide & \$0.11 \\
Not in the transcript (side calls) & \$68.83 \\
\hline
\hshead{Total} & \hsaccenthead{\$97.57} \\
\hline
\end{tabular}}
\end{hsblock}

The split follows hwde's own record of when each step ran, across ten
session transcripts. One session (\texttt{811c9ae8}) was still open when
this was recorded, so its \$0.34 above is its priced-token estimate, not a
final figure. About 29.1 million tokens ran on a model
(\texttt{claude-fable-5-1}) with no listed price; that gap between the
step-by-step transcripts and the session's own recorded total is folded into
the "side calls" row rather than split further.

\end{document}
